\documentclass[11pt,a4paper]{article}

% Pacotes de idioma e fontes
\usepackage[utf8]{inputenc}
\usepackage[T1]{fontenc}
\usepackage[portuguese]{babel}

% Pacotes Matemáticos e Gráficos
\usepackage{amsmath}
\usepackage{amsfonts}
\usepackage{amssymb}
\usepackage{tikz} 
\usepackage{pgfplots}
\pgfplotsset{compat=1.18}

% Pacotes de Layout
\usepackage[margin=1.5cm]{geometry}
\usepackage{paracol} % Permite colunas paralelas independentes
\usepackage{enumitem}

% Definição Global de Macros Gráficas
\newcommand{\match}[4]{
    \draw[line width=1.5pt, draw=gray!80] (#1,#2) -- (#3,#4);
    \fill[black!70] (#1,#2) circle (2pt);
}

% Configuração da linha divisória e larguras das colunas
\setlength{\columnsep}{1cm} 
\setlength{\columnseprule}{0.5pt}

\begin{document}

% Cabeçalho Centralizado
\begin{center}
    \textbf{\Large PEAC-2026 | MATEMÁTICA - EAD} \\[0.2cm]
    \textbf{\large AULA 16-17 - Função Exponencial e Função Logarítmica} \\[0.2cm]
    \textit{Questões de Vestibulares de anos anteriores (UFRGS).}
\end{center}
\vspace{0.5cm}

% Início do ambiente de colunas paralelas
\begin{paracol}{2}

% =========================================================
% QUESTÕES
% =========================================================

\switchcolumn[0]*
\textbf{1) UFRGS-2025} \\
Considere as seguintes afirmações sobre números racionais.

\begin{enumerate}[label=\Roman*., itemsep=0.2cm]
    \item Se $0 < a \le b$, então $b^a$ é sempre um número irracional.
    \item Se $0 < a \le b$, então $0 < a \le \dfrac{a+b}{2} \le b$.
    \item Se $0 < a < b < 1$, então $\log(a) < \log(b) < 0$.
\end{enumerate}

Quais estão corretas?

\begin{enumerate}[label=(\Alph*)]
    \item Apenas I.
    \item Apenas II.
    \item Apenas I e III.
    \item Apenas II e III.
    \item I, II e III.
\end{enumerate}
\vspace{0.4cm}

\switchcolumn[1]
\vspace{7cm} 

\switchcolumn[0]*
\textbf{2) UFRGS-2025} \\
A soma das raízes da equação $2^{x^2-3x-10} = 1$ é

\begin{enumerate}[label=(\Alph*)]
    \item $-3$.
    \item $-2$.
    \item $0$.
    \item $2$.
    \item $3$.
\end{enumerate}
\vspace{0.4cm}

\switchcolumn[1]
\vspace{6cm}

\switchcolumn[0]*
\textbf{3) UFRGS-2025} \\
A equação $\log(x) + \log(x - 1) = \log 6$ tem conjunto solução nos números reais. 

O conjunto solução da equação é

\begin{enumerate}[label=(\Alph*)]
    \item $\{-2\}$.
    \item $\{3\}$.
    \item $\{3, 2\}$.
    \item $\{3, -2\}$.
    \item $\{-3, -2\}$.
\end{enumerate}
\vspace{0.4cm}

\switchcolumn[1]
\vspace{6cm}

\switchcolumn[0]*
\textbf{4) UFRGS-2024} \\
Considere as seguintes afirmações sobre números e suas operações.

\begin{enumerate}[label=\Roman*., itemsep=0.2cm]
    \item A soma de dois números naturais consecutivos é sempre um número ímpar.
    \item A soma de dois números irracionais é sempre um número irracional.
    \item Se $a > 1$, então, para qualquer valor inteiro de $n$, $0 < \dfrac{1}{a^n} < 1$.
\end{enumerate}

Quais estão corretas?

\begin{enumerate}[label=(\Alph*)]
    \item Apenas I.
    \item Apenas III.
    \item Apenas I e II.
    \item Apenas II e III.
    \item I, II e III.
\end{enumerate}
\vspace{0.4cm}

\switchcolumn[1]
\vspace{7cm} 

\switchcolumn[0]*
\textbf{5) UFRGS-2024} \\
O valor da expressão \\
$\log(2024) - \log(202{,}4) + \log(20{,}24) - \log(2{,}024)$ é

\begin{enumerate}[label=(\Alph*)]
    \item 4.
    \item 2.
    \item 0.
    \item $-2$.
    \item $-4$.
\end{enumerate}
\vspace{0.4cm}

\switchcolumn[1]
\vspace{6cm}

\switchcolumn[0]*
\textbf{6) UFRGS-2023} \\
Considere as seguintes afirmações sobre números e suas operações.

\begin{enumerate}[label=\Roman*., itemsep=0.3cm]
    \item $1 + \dfrac{1}{2} + \dfrac{1}{4} + \dfrac{1}{8} \dots > 2 + \dfrac{2}{3} + \dfrac{2}{9} + \dfrac{2}{27} \dots \, .$
    \item $\sqrt{\sqrt{7} + 10} < \sqrt{\sqrt{7}} + \sqrt{10} \, .$
    \item $6 \cdot 5^{10} < 5 \cdot 6^{10} \, .$
\end{enumerate}

Quais estão corretas?

\begin{enumerate}[label=(\Alph*)]
    \item Apenas I.
    \item Apenas II.
    \item Apenas I e III.
    \item Apenas II e III.
    \item I, II e III.
\end{enumerate}
\vspace{0.4cm}

\switchcolumn[1]
\vspace{8cm} 

\switchcolumn[0]*
\textbf{7) UFRGS-2023} \\
Se $a$ e $b$ são as raízes da equação $x^2 + 2x - 15 = 0$, então o valor de $(ab)^{a+b}$ é

\begin{enumerate}[label=(\Alph*), itemsep=0.2cm]
    \item $-225$.
    \item $-\dfrac{1}{225}$.
    \item $-30$.
    \item $\dfrac{1}{225}$.
    \item $225$.
\end{enumerate}
\vspace{0.4cm}

\switchcolumn[1]
\vspace{7cm}

\switchcolumn[0]*
\textbf{8) UFRGS-2023} \\
A figura abaixo mostra o início de uma sequência infinita de quadrados. A medida dos lados dos quadrados $Q_1$, $Q_2$ e $Q_3$ são, respectivamente, $\log(2)$, $\log(\sqrt{2})$ e $\log(\sqrt[4]{2})$.

\begin{center}
\begin{tikzpicture}[scale=1.2]
    % Q1
    \draw[thick] (0,0) rectangle (2,2);
    \node[below] at (1,0) {$Q_1$};
    
    % Q2
    \draw[thick] (2.8,0) rectangle (3.8,1);
    \node[below] at (3.3,0) {$Q_2$};
    
    % Q3
    \draw[thick] (4.5,0) rectangle (5,0.5);
    \node[below] at (4.75,0) {$Q_3$};
    
    % Reticências
    \node[right] at (5.3,0.25) {...};
\end{tikzpicture}
\end{center}

A soma das áreas dessa sequência infinita de quadrados é

\begin{enumerate}[label=(\Alph*), itemsep=0.2cm]
    \item $\dfrac{1}{3} \cdot [\log(2)]^2$.
    \item $\dfrac{4}{3} \cdot [\log(2)]^2$.
    \item $\dfrac{2}{3} \cdot [\log(2)]^2$.
    \item $\log(2 + \sqrt{2} + \sqrt[4]{2})$.
    \item $\log(2 \cdot \sqrt{2} \cdot \sqrt[4]{2})$.
\end{enumerate}
\vspace{0.4cm}

\switchcolumn[1]
\vspace{9.5cm}

\switchcolumn[0]*
\textbf{9) UFRGS-2022} \\
O valor de $\log(2^2) + \log(2^3) + \log(2^4) + \dots + \log(2^{50})$ é

\begin{enumerate}[label=(\Alph*), itemsep=0.2cm]
    \item $\log(2^{1247})$.
    \item $\log(2^{1274})$.
    \item $\log(2^{1472})$.
    \item $\log(2^{59})$.
    \item $\log(8^{59})$.
\end{enumerate}
\vspace{0.4cm}

\switchcolumn[1]
\vspace{6cm}

\switchcolumn[0]*
\textbf{10) UFRGS-2021} \\
O valor de $\log\left(1 - \dfrac{1}{2}\right) + \log\left(1 - \dfrac{1}{3}\right) + \dots + \log\left(1 - \dfrac{1}{1000}\right)$ é

\begin{enumerate}[label=(\Alph*)]
    \item $-3$.
    \item $-2$.
    \item $-1$.
    \item $0$.
    \item $1$.
\end{enumerate}
\vspace{0.4cm}

\switchcolumn[1]
\vspace{6cm}

\switchcolumn[0]*
\textbf{11) UFRGS-2021} \\
Se $\log_2 x + (\log_2 x)^2 = 12$, então o valor de $x$ é

\begin{enumerate}[label=(\Alph*), itemsep=0.2cm]
    \item $8 \text{ ou } \dfrac{1}{16}$.
    \item $\dfrac{1}{8} \text{ ou } 16$.
    \item $-4 \text{ ou } 3$.
    \item $4 \text{ ou } -3$.
    \item $2^{12}$.
\end{enumerate}
\vspace{0.4cm}

\switchcolumn[1]
\vspace{6cm}

\switchcolumn[0]*
\textbf{12) UFRGS-2019} \\
A concentração de alguns medicamentos no organismo está relacionada com a meia-vida, ou seja, o tempo necessário para que a quantidade inicial do medicamento no organismo seja reduzida pela metade.

Considere que a meia-vida de determinado medicamento é de 6 horas. Sabendo que um paciente ingeriu 120 mg desse medicamento às 10 horas, assinale a alternativa que representa a melhor aproximação para a concentração desse medicamento, no organismo desse paciente, às 16 horas do dia seguinte.

\begin{enumerate}[label=(\Alph*)]
    \item 2,75 mg.
    \item 3 mg.
    \item 3,75 mg.
    \item 4 mg.
    \item 4,25 mg.
\end{enumerate}
\vspace{0.4cm}

\switchcolumn[1]
\vspace{8.5cm}

\switchcolumn[0]*
\textbf{13) UFRGS-2019} \\
Se $\log 2 = x$ e $\log 3 = y$, então $\log 288$ é

\begin{enumerate}[label=(\Alph*)]
    \item $2x + 5y$.
    \item $5x + 2y$.
    \item $10xy$.
    \item $x^2 + y^2$.
    \item $x^2 - y^2$.
\end{enumerate}
\vspace{0.4cm}

\switchcolumn[1]
\vspace{6cm}

\switchcolumn[0]*
\textbf{14) UFRGS-2018} \\
Sendo $a$ e $b$ números reais positivos quaisquer, considere as afirmações abaixo.

\begin{enumerate}[label=\Roman*., itemsep=0.2cm]
    \item Se $a > b$ então $a^x > b^x$, para qualquer $x \in \mathbf{R}$.
    \item Se $a > b > 1$, então $\dfrac{1}{a} < \dfrac{2}{a+b} < \dfrac{1}{b}$.
    \item Se $a > b$, então $\sqrt{a} > \sqrt{b}$.
\end{enumerate}

Quais estão corretas?

\begin{enumerate}[label=(\Alph*)]
    \item Apenas I.
    \item Apenas II.
    \item Apenas I e III.
    \item Apenas II e III.
    \item I, II e III.
\end{enumerate}
\vspace{0.4cm}

\switchcolumn[1]
\vspace{7.5cm}

\switchcolumn[0]*
\textbf{15) UFRGS-2018} \\
Dadas as funções reais de variável real $f$ e $g$, definidas por $f(x) = -\log_2(x)$ e $g(x) = x^2 - 4$, pode-se afirmar que $f(x) = g(x)$ é verdadeiro para um valor de $x$ localizado no intervalo

\begin{enumerate}[label=(\Alph*)]
    \item $[0; 1]$.
    \item $[1; 2]$.
    \item $[2; 3]$.
    \item $[3; 4]$.
    \item $[4; 5]$.
\end{enumerate}
\vspace{0.4cm}

\switchcolumn[1]
\vspace{6cm}

\switchcolumn[0]*
\textbf{16) UFRGS-2018} \\
Considere a função real de variável real $f(x) = 2^{x-1}$. Com relação à $f(x)$, é correto afirmar que

\begin{enumerate}[label=(\Alph*), itemsep=0.2cm]
    \item se $x < 1$, então $f(x) < 0$.
    \item se $x \ge 1$, então $f(x) \le 1$.
    \item a função $f(x)$ é decrescente para $x < 0$ e crescente para $x \ge 0$.
    \item os valores das imagens de $f(x): A \rightarrow \mathbf{R}$, em que $A = \{x \in \mathbf{N} / x \ge 0\}$, formam uma progressão aritmética.
    \item os valores das imagens de $f(x): A \rightarrow \mathbf{R}$, em que $A = \{x \in \mathbf{N} / x \ge 0\}$, formam uma progressão geométrica.
\end{enumerate}
\vspace{0.4cm}

\switchcolumn[1]
\vspace{8.5cm}

\switchcolumn[0]*
\textbf{17) UFRGS-2018} \\
O valor de \\
$E = \log\left(\dfrac{1}{2}\right) + \log\left(\dfrac{2}{3}\right) + \dots + \log\left(\dfrac{999}{1000}\right)$ é

\begin{enumerate}[label=(\Alph*)]
    \item $-3$.
    \item $-2$.
    \item $-1$.
    \item $0$.
    \item $1$.
\end{enumerate}
\vspace{0.4cm}

\switchcolumn[1]
\vspace{6cm}

\switchcolumn[0]*
\textbf{18) UFRGS-2017} \\
Considere a função real $f$ definida por $f(x) = 2^{-x}$. O valor da expressão $S = f(0) + f(1) + f(2) + \dots + f(100)$ é

\begin{enumerate}[label=(\Alph*), itemsep=0.3cm]
    \item $S = 2 - 2^{-101}$.
    \item $S = 2^{50} + 2^{-50}$.
    \item $S = 2 + 2^{-101}$.
    \item $S = 2 + 2^{-100}$.
    \item $S = 2 - 2^{-100}$.
\end{enumerate}
\vspace{0.4cm}

\switchcolumn[1]
\vspace{8cm}

\switchcolumn[0]*
\textbf{19) UFRGS-2017} \\
Se $\log_3 x + \log_9 x = 1$, então o valor de $x$ é

\begin{enumerate}[label=(\Alph*), itemsep=0.3cm]
    \item $\sqrt[3]{2}$.
    \item $\sqrt{2}$.
    \item $\sqrt[3]{3}$.
    \item $\sqrt{3}$.
    \item $\sqrt[3]{9}$.
\end{enumerate}
\vspace{0.4cm}

\switchcolumn[1]
\vspace{8cm}

\switchcolumn[0]*
\textbf{20) UFRGS-2017} \\
Leia o texto abaixo, sobre terremotos.

Magnitude é uma medida quantitativa do tamanho do terremoto. Ela está relacionada com a energia sísmica liberada no foco e também com a amplitude das ondas registradas pelos sismógrafos. Para cobrir todos os tamanhos de terremotos, desde os microtremores de magnitudes negativas até os grandes terremotos com magnitudes superiores a 8.0, foi idealizada uma escala logarítmica, sem limites. No entanto, a própria natureza impõe um limite superior a esta escala, já que ela está condicionada ao próprio limite de resistência das rochas da crosta terrestre. Magnitude e energia podem ser relacionadas pela fórmula descrita por Gutenberg e Richter em 1935: $\log(E) = 11{,}8 + 1{,}5M$ onde: $E =$ energia liberada em \textit{Erg} ; $M =$ magnitude do terremoto.

\begin{flushright}
\footnotesize
\textbf{Disponível em: <http://www.iag.usp.br/siae98/terremoto/terremotos.htm>.} \\
\textbf{Acesso em: 20 set. 2017.}
\end{flushright}
\vspace{0.2cm}

Sabendo que o terremoto que atingiu o México em setembro de 2017 teve magnitude 8,2, assinale a alternativa que representa a melhor aproximação para a energia liberada por esse terremoto, em \textit{Erg}.

\begin{enumerate}[label=(\Alph*), itemsep=0.2cm]
    \item $13{,}3$
    \item $20$
    \item $24$
    \item $10^{24}$
    \item $10^{28}$
\end{enumerate}
\vspace{0.4cm}

\switchcolumn[1]
\vspace{12cm}

\switchcolumn[0]*
\textbf{21) UFRGS-2016} \\
Se $\log_5 x = 2$ e $\log_{10} y = 4$, então $\log_{20} \dfrac{y}{x}$ é

\begin{enumerate}[label=(\Alph*)]
    \item 2.
    \item 4.
    \item 6.
    \item 8.
    \item 10.
\end{enumerate}
\vspace{0.4cm}

\switchcolumn[1]
\vspace{6cm}

\switchcolumn[0]*
\textbf{22) UFRGS-2016} \\
No estudo de uma população de bactérias, identificou-se que o número $N$ de bactérias, $t$ horas após o início do estudo, é dado por $N(t) = 20 \cdot 2^{1{,}5t}$.

Nessas condições, em quanto tempo a população de mosquitos duplicou?

\begin{enumerate}[label=(\Alph*)]
    \item 15 min.
    \item 20 min.
    \item 30 min.
    \item 40 min.
    \item 45 min.
\end{enumerate}
\vspace{0.4cm}

\switchcolumn[1]
\vspace{7cm}

\switchcolumn[0]*
\textbf{23) UFRGS-2015} \\
Se $10^x = 20^y$, atribuindo $0{,}3$ para $\log 2$, então o valor de $\dfrac{x}{y}$ é

\begin{enumerate}[label=(\Alph*)]
    \item 0,3.
    \item 0,5.
    \item 0,7.
    \item 1.
    \item 1,3.
\end{enumerate}
\vspace{0.4cm}

\switchcolumn[1]
\vspace{6cm}

\switchcolumn[0]*
\textbf{24) UFRGS-2015} \\
Considere a função $f$ definida por $f(x) = 1 - 5 \cdot 0{,}7^x$ e representada em um sistema de coordenadas cartesianas.

Entre os gráficos abaixo, o que pode representar a função $f$ é

\vspace{0.4cm}
\begin{tabular}{@{}ll@{}}
(A) \begin{tikzpicture}[scale=0.35, baseline=-0.5ex]
    \draw[->, thick] (-2,0) -- (4,0) node[below] {\tiny $x$};
    \draw[->, thick] (0,-4) -- (0,3) node[left] {\tiny $y$};
    \node[below left] at (0,0) {\tiny $0$};
    \draw[domain=-1.2:3.5, smooth, variable=\x, thick] plot ({\x}, {1.5 - 2*exp(-1.2*\x)});
\end{tikzpicture} &
(B) \begin{tikzpicture}[scale=0.35, baseline=-0.5ex]
    \draw[->, thick] (-1,0) -- (4,0) node[below] {\tiny $x$};
    \draw[->, thick] (0,-1) -- (0,4) node[left] {\tiny $y$};
    \node[below left] at (0,0) {\tiny $0$};
    \draw[domain=-0.8:3.5, smooth, variable=\x, thick] plot ({\x}, {0.2 + 2.5*exp(-1.5*\x)});
\end{tikzpicture} \\[1cm]
(C) \begin{tikzpicture}[scale=0.35, baseline=-0.5ex]
    \draw[->, thick] (-3,0) -- (4,0) node[below] {\tiny $x$};
    \draw[->, thick] (0,-3) -- (0,3) node[left] {\tiny $y$};
    \node[below left] at (0,0) {\tiny $0$};
    \draw[domain=-1.5:3.5, smooth, variable=\x, thick] plot ({\x}, {-1.5 + 0.8*exp(-1.2*\x)});
\end{tikzpicture} &
(D) \begin{tikzpicture}[scale=0.35, baseline=-0.5ex]
    \draw[->, thick] (-3,0) -- (4,0) node[below] {\tiny $x$};
    \draw[->, thick] (0,-3) -- (0,3) node[left] {\tiny $y$};
    \node[below left] at (0,0) {\tiny $0$};
    \draw[domain=-1.5:3.5, smooth, variable=\x, thick] plot ({\x}, {1.5 - 0.4*exp(-1.2*\x)});
\end{tikzpicture} \\[1cm]
(E) \begin{tikzpicture}[scale=0.35, baseline=-0.5ex]
    \draw[->, thick] (-2,0) -- (4,0) node[below] {\tiny $x$};
    \draw[->, thick] (0,-3) -- (0,3) node[left] {\tiny $y$};
    \node[below left] at (0,0) {\tiny $0$};
    \draw[thick] plot[smooth, tension=0.7] coordinates {(-1.5, 3) (-0.5, 0) (0, -1) (1, -1.5) (3, -0.5) (4, -0.2)};
\end{tikzpicture} &
\end{tabular}
\vspace{0.4cm}

\switchcolumn[1]
\vspace{16cm}

\switchcolumn[0]*
\textbf{25) UFRGS-2014} \\
A expressão $(0{,}125)^{15}$ é equivalente a

\begin{enumerate}[label=(\Alph*)]
    \item $5^{45}$.
    \item $5^{-45}$.
    \item $2^{45}$.
    \item $2^{-45}$.
    \item $(-2)^{45}$.
\end{enumerate}
\vspace{0.4cm}

\switchcolumn[1]
\vspace{6cm}

\switchcolumn[0]*
\textbf{26) UFRGS-2014} \\
O algarismo das unidades de $9^{99} - 4^{44}$ é

\begin{enumerate}[label=(\Alph*)]
    \item 1.
    \item 2.
    \item 3.
    \item 4.
    \item 5.
\end{enumerate}
\vspace{0.4cm}

\switchcolumn[1]
\vspace{6cm}

\switchcolumn[0]*
\textbf{27) UFRGS-2014} \\
Por qual potência de $10$ deve ser multiplicado o número $10^{-3} \cdot 10^{-3} \cdot 10^{-3} \cdot 10^{-3}$ para que esse produto seja igual a $10$?

\begin{enumerate}[label=(\Alph*)]
    \item $10^9$.
    \item $10^{10}$.
    \item $10^{11}$.
    \item $10^{12}$.
    \item $10^{13}$.
\end{enumerate}
\vspace{0.4cm}

\switchcolumn[1]
\vspace{6cm}

\switchcolumn[0]*
\textbf{28) UFRGS-2014} \\
Atribuindo para $\log 2$ o valor $0{,}3$, então o valor de $100^{0{,}3}$ é

\begin{enumerate}[label=(\Alph*)]
    \item 3.
    \item 4.
    \item 8.
    \item 10.
    \item 33.
\end{enumerate}
\vspace{0.4cm}

\switchcolumn[1]
\vspace{6cm}

\switchcolumn[0]*
\textbf{29) UFRGS-2014} \\
O número $N$ de peixes em um lago pode ser estimado utilizando a função $N$, definida por $N(t) = 500 \cdot 1{,}02^t$, em que $t$ é o tempo medido em meses.

Pode-se, então, estimar que a população de peixes no lago, a cada mês,

\begin{enumerate}[label=(\Alph*)]
    \item cresce $0{,}2\%$.
    \item cresce $2\%$.
    \item cresce $20\%$.
    \item decresce $2\%$.
    \item decresce $20\%$.
\end{enumerate}
\vspace{0.4cm}

\switchcolumn[1]
\vspace{7cm}

\switchcolumn[0]*
\textbf{30) UFRGS-2012} \\
O algarismo das unidades da soma $44^{54} + 55^{45}$ é

\begin{enumerate}[label=(\Alph*)]
    \item 0.
    \item 1.
    \item 2.
    \item 3.
    \item 4.
    \item 5.
\end{enumerate}
\vspace{0.4cm}

\switchcolumn[1]
\vspace{6cm}

\switchcolumn[0]*
\textbf{31) UFRGS-2012} \\
Dez bactérias são cultivadas para uma experiência, e o número de bactérias dobra a cada 12 horas.

Tomando como aproximação para $\log 2$ o valor $0{,}3$, decorrida exatamente uma semana, o número de bactérias está entre

\begin{enumerate}[label=(\Alph*), itemsep=0.2cm]
    \item $10^{4{,}5} \text{ e } 10^5$.
    \item $10^5 \text{ e } 10^{5{,}5}$.
    \item $10^{5{,}5} \text{ e } 10^6$.
    \item $10^6 \text{ e } 10^{6{,}5}$.
    \item $10^{6{,}5} \text{ e } 10^7$.
\end{enumerate}
\vspace{0.4cm}

\switchcolumn[1]
\vspace{7cm}

\switchcolumn[0]*
\textbf{32) UFRGS-2011} \\
O número $\log_2 7$ está entre

\begin{enumerate}[label=(\Alph*)]
    \item $0 \text{ e } 1$.
    \item $1 \text{ e } 2$.
    \item $2 \text{ e } 3$.
    \item $3 \text{ e } 4$.
    \item $4 \text{ e } 5$.
\end{enumerate}
\vspace{0.4cm}

\switchcolumn[1]
\vspace{6cm}

\switchcolumn[0]*
\textbf{33) UFRGS-2011} \\
Considere a função $f$ tal que $f(x) = k + \left(\dfrac{5}{4}\right)^{2x-1}$, com $k > 0$.

Assinale a alternativa correspondente ao gráfico que pode representar a função $f$.

\vspace{0.4cm}
\begin{tabular}{@{}ll@{}}
(A) \begin{tikzpicture}[scale=0.35, baseline=-0.5ex]
    \draw[->, thick] (-3,0) -- (3,0) node[below] {\tiny $x$};
    \draw[->, thick] (0,-1) -- (0,5) node[left] {\tiny $y$};
    \node[below left] at (0,0) {\tiny $0$};
    \draw[gray!60, dashed] (-3,1.2) -- (3,1.2);
    \draw[domain=-3:2.2, smooth, variable=\x, thick] plot ({\x}, {1.2 + 0.5*exp(0.9*\x)});
\end{tikzpicture} &
(B) \begin{tikzpicture}[scale=0.35, baseline=-0.5ex]
    \draw[->, thick] (-3,0) -- (3,0) node[below] {\tiny $x$};
    \draw[->, thick] (0,-2) -- (0,5) node[left] {\tiny $y$};
    \node[below left] at (0,0) {\tiny $0$};
    \draw[gray!60, dashed] (-3,-1) -- (3,-1);
    \draw[domain=-3:2.3, smooth, variable=\x, thick] plot ({\x}, {-1 + 0.5*exp(1.1*\x)});
\end{tikzpicture} \\[1cm]
(C) \begin{tikzpicture}[scale=0.35, baseline=-0.5ex]
    \draw[->, thick] (-3,0) -- (3,0) node[below] {\tiny $x$};
    \draw[->, thick] (0,-2) -- (0,5) node[left] {\tiny $y$};
    \node[below left] at (0,0) {\tiny $0$};
    \draw[gray!60, dashed] (-3,-1) -- (3,-1);
    \draw[domain=-2.3:3, smooth, variable=\x, thick] plot ({\x}, {-1 + 0.5*exp(-1.1*\x)});
\end{tikzpicture} &
(D) \begin{tikzpicture}[scale=0.35, baseline=-0.5ex]
    \draw[->, thick] (-3,0) -- (3,0) node[below] {\tiny $x$};
    \draw[->, thick] (0,-1) -- (0,5) node[left] {\tiny $y$};
    \node[below left] at (0,0) {\tiny $0$};
    \draw[gray!60, dashed] (-3,1.2) -- (3,1.2);
    \draw[domain=-2.2:3, smooth, variable=\x, thick] plot ({\x}, {1.2 + 0.5*exp(-0.9*\x)});
\end{tikzpicture} \\[1cm]
(E) \begin{tikzpicture}[scale=0.35, baseline=-0.5ex]
    \draw[->, thick] (-3,0) -- (3,0) node[below] {\tiny $x$};
    \draw[->, thick] (0,-4) -- (0,3) node[left] {\tiny $y$};
    \node[below left] at (0,0) {\tiny $0$};
    \draw[gray!60, dashed] (-3,1) -- (3,1);
    \draw[domain=-3:1.8, smooth, variable=\x, thick] plot ({\x}, {1 - 0.4*exp(1.2*\x)});
\end{tikzpicture} &
\end{tabular}
\vspace{0.4cm}

\switchcolumn[1]
\vspace{16cm}

\switchcolumn[0]*
\textbf{34) UFRGS-2010} \\
Aproximando $\log 2$ por $0{,}301$, verificamos que o número $16^{10}$ está entre

\begin{enumerate}[label=(\Alph*)]
    \item $10^9 \text{ e } 10^{10}$.
    \item $10^{10} \text{ e } 10^{11}$.
    \item $10^{11} \text{ e } 10^{12}$.
    \item $10^{12} \text{ e } 10^{13}$.
    \item $10^{13} \text{ e } 10^{14}$.
\end{enumerate}
\vspace{0.4cm}

\switchcolumn[1]
\vspace{6cm}

\end{paracol}

\newpage
% ---------------------------------------------------------
% GABARITO
% ---------------------------------------------------------
\begin{center}
    \rule{0.6\textwidth}{0.5pt} \\[0.3cm]
    \textbf{\large GABARITO} \\[0.3cm]
    \begin{tabular}{|c|c|c|c|c|c|c|c|}
        \hline
        \textbf{1} & D & \textbf{2} & E & \textbf{3} & B & \textbf{4} & A \\ \hline
        \textbf{5} & B & \textbf{6} & D & \textbf{7} & D & \textbf{8} & B \\ \hline
        \textbf{9} & B & \textbf{10} & A & \textbf{11} & A & \textbf{12} & C \\ \hline
        \textbf{13} & B & \textbf{14} & D & \textbf{15} & B & \textbf{16} & E \\ \hline
        \textbf{17} & A & \textbf{18} & E & \textbf{19} & E & \textbf{20} & D \\ \hline
        \textbf{21} & A & \textbf{22} & D & \textbf{23} & E & \textbf{24} & A \\ \hline
        \textbf{25} & D & \textbf{26} & C & \textbf{27} & E & \textbf{28} & B \\ \hline
        \textbf{29} & B & \textbf{30} & B & \textbf{31} & B & \textbf{32} & C \\ \hline
        \textbf{33} & A & \textbf{34} & D & \multicolumn{4}{c|}{} \\ \hline
    \end{tabular} \\[0.3cm]
    \rule{0.6\textwidth}{0.5pt}
\end{center}

\end{document}